\chapter{Source Code Reference}

The whole program is one file, \texttt{three-dimension-maze-game/main.cpp}, of about 458 lines (excluding comments/whitespace). It uses only the C++ standard library (C++20) and SDL2 (window, renderer, texture, input, timers). No BGI, no other external dependencies.

\section{Includes and Constants}
\begin{minted}{cpp}
#include <SDL2/SDL.h>
#include <bits/stdc++.h>           // vector, algorithm, cmath, etc.
#define PI 3.14159265358979323846
using namespace std;

const int       SCREEN_WIDTH    = 640;
const int       SCREEN_HEIGHT   = 480;
const double    EYE_HEIGHT      = 0.5;     // camera height (half a wall)
const double    FOV_DEG         = 75.0;    // field of view, degrees
const double    MOVE_SPD        = 2.0;     // units per second (ground)
const double    TURN_SPD        = 2.5;     // radians per second
const double    CAPS_MULTIPLIER = 2.5;     // sprint multiplier (Shift/CapsLock)
const double    WALL_H          = 1.0;     // wall height (matches EYE_HEIGHT)
const double    PLAYER_R        = 0.28;    // collision radius (corridor = 1.0)
\end{minted}

\section{Vector Math --- \texttt{vec3}}
The \texttt{vec3} class and its associated free functions (operators, \texttt{dot}, \texttt{cross}, \texttt{unit\_vector}) are adapted from \textit{Ray Tracing in One Weekend} by Peter Shirley \cite{shirley2016ray}. We use a minimal subset sufficient for this project: a 3-component double array with x/y/z accessors, basic arithmetic operators, dot/cross products, and vector normalization.

\section{Ray and Hit Record}
\begin{minted}{cpp}
struct ray {
    vec3 origin, direction;
    ray(vec3 o, vec3 d) : origin(o), direction(d) {}
    vec3 at(double t) const { return origin + direction*t; }
};

struct hit_record {
    double ray_t = 1e30;            // distance along ray (huge = no hit)
    vec3 hit_point, normal, material_color;
};
\end{minted}

\section{Core Pipeline Functions}

\subsection{\texttt{is\_wall} (inline, static)}
\begin{minted}{cpp}
static bool is_wall(int cx, int cy) {
    if (cx < 0 || cy < 0 || cx >= MAZE_W || cy >= MAZE_H) return true;  // outside = solid
    return MAP[cy][cx] == '#';
}
\end{minted}


\subsection{\texttt{build\_exit\_dist} (static)}
BFS over floor cells from the exit cell. Queue implemented as growing \texttt{vector} (index loop). \texttt{max\_dist} recorded for temperature normalization.

\subsection{\texttt{floor\_color\_at} (returns green on exit cell, checkerboard otherwise)}
Shown fully in Section 2 (green: $\texttt{vec3(0.20, 0.85, 0.30)}$; light gray: $\texttt{vec3(0.60, 0.58, 0.55)}$; dark gray: $\texttt{vec3(0.30, 0.30, 0.33)}$).

\subsection{\texttt{temperature (static)}
Maps BFS distance $d$ to $[-1, 1]$; $-1$ for unreachable cells (set to $\text{max\_dist}$).

\subsection{\texttt{cast\_ray (ray) $\to$ hit\_record}
Shown fully in Section 2. Tests floor plane ($z=0$), ceiling plane ($z=WALL\_H$), then 2D DDA walk (bounded 256 iterations, early exit on first wall hit, altitude check at entry). Uses \texttt{unit\_vector} only for direction normalization; ray direction from render is raw.

\subsection{\texttt{shade (hit\_record, eye) $\to$ vec3}
Shown fully in Section 2. Headlamp Lambertian with warm color $(1.7, 1.45, 1.10)$, ambient $(0.14, 0.14, 0.18)$, attenuation $1/(0.9 + 0.30 d^2)$, gamma $\sqrt{\cdot}$.

\subsection{\texttt{collides (pos) $\to$ bool}
Shown fully in Section 2. 2D circle (radius PLAYER\_R) vs axis-aligned squares. Checks cells in $[\text{floor}(p_x \pm R), \text{floor}(p_x \pm R)]$; clamp to find closest point; squared distance check.

\subsection{\texttt{render (buffer, eye, yaw)}
Shown fully in Section 2. Computes camera basis ($\text{forward} = (\cos\psi, \sin\psi, 0)$, $\text{right} = \text{cross}(\text{forward}, (0,0,1))$), then per-pixel $u,v$, constructs ray, calls \texttt{cast\_ray} + \texttt{shade}, packs ARGB32.

\subsection{\texttt{update (pos, yaw, dt, keys)}
Shown fully in Section 2. Handles yaw (left/right), walk (W/S/Up/Down), strafe (A/D), sprint (Shift/CapsLock), axis-separated commit with \texttt{collides}.

\section{Main Loop}
\label{sec:mainloop}
\begin{minted}{cpp}
int32_t main() {
    vector<Uint32> buffer(SCREEN_WIDTH * SCREEN_HEIGHT);
    Uint32 pitch = SCREEN_WIDTH * sizeof(Uint32);

    // SDL init
    SDL_SetMainReady();
    SDL_Init(SDL_INIT_VIDEO);
    SDL_Window*   window   = SDL_CreateWindow("Maze Runner 3D", ...);
    SDL_Renderer* renderer = SDL_CreateRenderer(window, -1, SDL_RENDERER_ACCELERATED | SDL_RENDERER_PRESENTVSYNC);
    SDL_Texture*  texture  = SDL_CreateTexture(renderer, SDL_PIXELFORMAT_ARGB8888, SDL_TEXTUREACCESS_STREAMING, SCREEN_WIDTH, SCREEN_HEIGHT);

    // Find S and E in MAP, build BFS
    for (int z = 0; z < MAZE_H; z++)
        for (int x = 0; x < MAZE_W; x++) {
            if (MAP[z][x] == 'S') { start_x = x; start_z = z; }
            if (MAP[z][x] == 'E') { exit_x  = x; exit_z  = z; }
        }
    build_exit_dist();

    vec3 pos(start_x + 0.5, start_z + 0.5, EYE_HEIGHT); // spawn center of S cell, EYE_HEIGHT
    double yaw = 0.0;                                    // facing +x

    SDL_Event event;
    bool quit = false, won = false;
    Uint32 last = SDL_GetTicks();

    while (!quit) {
        while (SDL_PollEvent(&event)) {
        if (event.type == SDL_QUIT || (event.type == SDL_KEYDOWN && event.key.keysym.sym == SDLK_ESCAPE))
            quit = true;
        }

        Uint32 now = SDL_GetTicks();
        double dt = (now - last) / 1000.0;
        last = now;
        if (dt > 0.05) dt = 0.05;           // clamp large stalls
        update(pos, yaw, dt, SDL_GetKeyboardState(NULL));

        if ((int)floor(pos.x()) == exit_x && (int)floor(pos.y()) == exit_z) {
            won = true; quit = true;
        }

        // Hot/cold score in title bar
        double temp = temperature(pos);
        const char* feel = temp > 0.5 ? "HOT" : temp > 0.0 ? "WARM" : temp > -0.5 ? "COOL" : "COLD";
        char title[64];
        snprintf(title, sizeof title, "Maze Runner 3D  |  Score %+.2f  %s", temp, feel);
        SDL_SetWindowTitle(window, title);

        render(buffer, vec3(pos.x(), pos.y(), EYE_HEIGHT), yaw);
        SDL_UpdateTexture(texture, NULL, buffer.data(), pitch);
        SDL_RenderCopy(renderer, texture, NULL, NULL);

        // Hot/cold bar (drawn on renderer, on top of texture)
        double f = (temp + 1.0) / 2.0;
        SDL_SetRenderDrawColor(renderer, 20, 24, 30, 255);
        SDL_Rect bg{16, 12, SCREEN_WIDTH - 32, 10}; SDL_RenderFillRect(renderer, &bg);
        SDL_SetRenderDrawColor(renderer, (Uint8)(60 + 195*f), (Uint8)(120 - 60*f), (Uint8)(255 - 215*f), 255);
        SDL_Rect bar{16, 12, (int)((SCREEN_WIDTH - 32) * f), 10}; SDL_RenderFillRect(renderer, &bar);

        SDL_RenderPresent(renderer);
        Uint32 spent = SDL_GetTicks() - now;
        if (spent < 15) SDL_Delay(15 - spent);
    }

    printf("%s after %.1f seconds.\n", won ? "Escaped the maze" : "Quit", (SDL_GetTicks() - last)/1000.0);
    // Cleanup
    SDL_DestroyTexture(texture); SDL_DestroyRenderer(renderer); SDL_DestroyWindow(window); SDL_Quit();
    return 0;
}
\end{minted}

\section{Build and Run}
\begin{verbatim}
cd three-dimension-maze-game
./compile.sh
# or manually:
g++ main.cpp -o main -I.. -std=c++20 -g -O3 -Wall -Wextra -Wpedantic -Wshadow $(sdl2-config --cflags --libs)
./main
\end{verbatim
The program requires SDL2 development libraries. On Debian/Ubuntu: \texttt{sudo apt install libsdl2-dev}. The hot/cold temperature score appears in the title bar and as a horizontal bar at the top of the window (blue = cold, red = hot). Find the green floor tile (exit) and step on it to win.
\end{parameter>
